\chapter[The winners, part 1: labels]{The winners, part 1: labels from AI agents}
\label{ch:winlabels}

\begin{tldr}
\begin{itemize}
\item First place: \textbf{Team Kangaroo} (Kaggle users tascj and sakaku), \textbf{0.969 public / 0.960 private}.
\item They treated the competition ``mainly as a labelling competition''.
\item For every training study an \textbf{AI agent} (a frontier model with tools) read the report \textbf{and looked at
      the images}, wrote down its evidence, then gave a probability per finding.
\item Their labels agreed with the doctors at \textbf{0.936--0.943} AUC (ours: 0.879). The same image model got
      +0.013 public from the labels alone.
\end{itemize}
\end{tldr}

\section{The big picture}

\begin{figure}[H]
\centering
\begin{tikzpicture}[x=1mm,y=1mm,
  src/.style={draw=grid, fill=white, rounded corners=1.5mm, align=center, font=\scriptsize, minimum height=11mm, text width=24mm},
  ag/.style={draw=them, fill=whyfill, rounded corners=1.5mm, align=center, font=\scriptsize, minimum height=11mm, text width=30mm},
  outb/.style={draw=them, fill=white, line width=0.8pt, rounded corners=1.5mm, align=center, font=\scriptsize, minimum height=11mm, text width=30mm},
  ar/.style={-{Stealth[length=1.8mm]}, line width=0.7pt, muted}]
  \node[src] (r) at (0,10) {\faFile\ report};
  \node[src] (i) at (0,-10) {\faImage\ MRI series};
  \node[ag] (a1) at (42,12) {\faRobot\ agent: 12 probabilities\\ (one session per study)};
  \node[ag] (a2) at (42,-12) {\faRobot\ agents: boxes around\\ anatomy and lesions};
  \node[outb] (o1) at (86,12) {\textbf{soft labels}\\ 4,349 $\times$ 12};
  \node[outb] (o2) at (86,-12) {\textbf{boxes}\\ (training only)};
  \node[draw=accent, fill=keyfill, rounded corners=2mm, align=center, font=\scriptsize, minimum height=30mm, text width=24mm]
     (m) at (125,0) {\faBrain\\ image classifier\\[1mm] at test time:\\ \textbf{images only}};
  \draw[ar] (r) -- (a1); \draw[ar] (i) -- (a1); \draw[ar] (i) -- (a2);
  \draw[ar] (a1) -- (o1); \draw[ar] (a2) -- (o2); \draw[ar] (o1) -- (m); \draw[ar] (o2) -- (m);
\end{tikzpicture}
\caption{Two annotation paths feed one image model. The agents are used only to make training data; the
submission runs the image model alone, offline. This chapter covers the upper path, Chapter~\ref{ch:winmodels} the lower.}
\label{fig:winoverview}
\end{figure}

\section{An agent, not a prompt}

Our labeller (Chapter~\ref{ch:labels}) was a \emph{prompt}: one fixed input, one answer. An \term{agent} is a
language model running in a loop with tools: it looks at something, decides what it still needs to know, uses a tool
to get it, and only answers when it is done.

\begin{figure}[H]
\centering
\begin{tikzpicture}[x=1mm,y=1mm,
  st/.style={draw=them, fill=whyfill, rounded corners=2mm, minimum width=36mm, minimum height=11mm, align=center, font=\scriptsize},
  ar/.style={-{Stealth[length=2mm]}, line width=0.8pt, them}]
  \node[st] (a) at (0,0) {\textbf{read}\\ instructions, report,\\ overview of each series};
  \node[st] (b) at (50,12) {\textbf{decide}\\ ``is the ACL intact?\\ check sagittal 14--17''};
  \node[st] (c) at (100,0) {\textbf{use a tool}\\ view those slices,\\ cropped and zoomed};
  \node[st] (d) at (50,-12) {\textbf{look}\\ at the returned image\\ sheet};
  \node[st, draw=wordline, fill=wordfill] (e) at (50,-34) {\textbf{answer}\\ evidence + 12 probabilities};
  \draw[ar] (a) -- (b); \draw[ar] (b) -- (c); \draw[ar] (c) -- (d); \draw[ar] (d) -- (b);
  \draw[ar, wordline] (d) -- (e) node[midway, right, font=\scriptsize] {when sure enough};
  \node[font=\scriptsize, text=muted, align=center] at (100,-16) {up to \ca 10 tool calls;\\ median 2.5 minutes\\ and 15 images per study};
\end{tikzpicture}
\caption{One labelling session, for one study. 4,349 such sessions labelled the training set.}
\label{fig:agentloop}
\end{figure}

\begin{wordcheck}
\term{Agent}: a language model that can call tools (run a script, view an image) in a loop. \term{Frontier model}:
the strongest hosted models of the time; here OpenAI's GPT-6 Astra and GPT-6.1 Sol, driven through command-line
agent tools (Codex, Claude Code and others). \term{Non-interactive session}: the agent runs alone, start to finish,
and returns an answer in a fixed format.
\end{wordcheck}

\section{The agent's workspace}

Each study got its own folder, prepared before the agent started:

\begin{figure}[H]
\centering
\begin{tikzpicture}[x=1mm,y=1mm, every node/.style={font=\scriptsize, anchor=west}]
  \node at (0,0) {\faFolder\ \code{<workspace>/}};
  \draw[muted] (2,-2) -- (2,-31);
  \foreach \y/\t/\c in {-5/{\faFile\ \code{AGENTS.md}}/{task, the 12 definitions and thresholds, series table, \textbf{the full report}, answer format},
                     -11/{\faFolder\ \code{overviews/}}/{one overview sheet per fluid series (up to 48 slices, 6 per row)},
                     -17/{\faFolder\ \code{series/}}/{every slice as a PNG, DICOM geometry, which series has which role},
                     -23/{\faFolder\ \code{tools/}}/{\code{view\_slices}: 1--4 slices, cropped, zoomed, adjustable brightness\\
                         \code{compare\_location}: where does this point lie in another series?},
                     -31/{\faFolder\ \code{scratch/views/}}/{the only writable folder, for the viewer's output}} {
    \draw[muted] (2,\y) -- (5,\y);
    \node at (5,\y) {\t};
    \node[text=muted, align=left] at (40,\y) {\c};
  }
\end{tikzpicture}
\caption{The workspace (after tascj's write-up). The agent sees one study only; switching the underlying model was
one line of configuration, so models and instructions could be compared fairly.}
\label{fig:workspace}
\end{figure}

\section{Evidence first, then a number}

The required answer format made the agent write its evidence \emph{before} its probability, per finding. In the style
of their format, for our Study K:

\begin{center}
\begin{tcolorbox}[enhanced, width=0.86\linewidth, colback=white, colframe=grid, boxrule=0.5pt, arc=1mm,
   fontupper=\scriptsize\ttfamily, left=3mm]
"acl": \{\\
\hspace*{3mm}"report\_evidence": ["Complete rupture of the anterior cruciate ligament."],\\
\hspace*{3mm}"image\_evidence": ["SAG\_FLUID slices 14-17: no continuous ACL fibres between femur\\
\hspace*{3mm}\hspace*{22mm}and tibia; fluid in the notch. Structural sagittal agrees."],\\
\hspace*{3mm}"present\_probability": 0.97\\
\}
\end{tcolorbox}
\end{center}

\begin{why}
Writing the evidence first forces the model to ``look before it judges'', and it makes every label checkable: a human
can see which slices the agent looked at and whether the claim came from the report or from the image.
\end{why}

\section{How good were the labels?}

\begin{figure}[H]
\centering
\begin{tikzpicture}
\begin{axis}[xbar, width=10.5cm, height=7.4cm, bar width=3mm, xmin=0.78, xmax=0.96,
    symbolic y coords={trained,astra2,astra1,hulu,qwen,clin,v4,v3},
    ytick={trained,astra2,astra1,hulu,qwen,clin,v4,v3}, yticklabels={{agent + trained guide (in-sample)}, {agent: Astra v2}, {agent: Astra v1},
      {Hulu-Med, report + 21 slices}, {Qwen Next, report + 21 slices}, {ClinFusion, report + 21 slices},
      {\textbf{our v4}, report only}, {\textbf{our v3}, report only}},
    y dir=reverse, y tick label style={font=\scriptsize}, xmajorgrids, xlabel={agreement with the 58 gold studies (macro AUC)},
    xlabel style={font=\footnotesize}, nodes near coords, nodes near coords style={font=\scriptsize, anchor=west},
    point meta=x, every node near coord/.append style={/pgf/number format/fixed, /pgf/number format/precision=3,
    /pgf/number format/fixed zerofill}, enlarge y limits=0.08]
  \addplot[barthem, bar shift=0pt] coordinates {(0.943,trained) (0.938,astra2) (0.936,astra1)};
  \addplot[fill=them!45, draw=white, bar shift=0pt] coordinates {(0.909,hulu) (0.896,qwen) (0.874,clin)};
  \addplot[barus, bar shift=0pt] coordinates {(0.879,v4) (0.809,v3)};
\end{axis}
\end{tikzpicture}
\caption{Label quality, same ruler for everybody. Single-pass readers that saw a fixed packet of 21 slices (7 per
plane) already matched or beat our report-only labels; agents that chose what to look at reached 0.94. The winners'
numbers are slightly optimistic: the gold studies were also used to develop their instructions.}
\label{fig:labelquality}
\end{figure}

\begin{innumbers}
Same image model (CoAtNet, 224\,px), only the labels swapped (earlier image-aware labels $\rightarrow$ agent labels):
\textbf{0.922 $\rightarrow$ 0.935} on the public leaderboard. For comparison, our own label upgrade (v3 $\rightarrow$
v4) gave +0.016.
\end{innumbers}

\section{Judging labels by more than AUC}

AUC only checks the order. For \emph{training} labels the values matter too (Chapter~\ref{ch:auc}), so they also
counted two things:
\begin{itemize}
\item \term{Log loss}: the average BCE of the labels against the gold answers. Lower = better calibrated.
\item \term{Confident errors}: labels that were confidently wrong (a high probability where the doctors said no, or
      the reverse), out of $58 \times 12 = 696$.
\end{itemize}

\begin{center}\small
\begin{tabular}{llccc}
\toprule
Model & Instructions & AUC & log loss & confident errors \\
\midrule
GPT-6 Luna & tascj's & 0.904 & 0.555 & 92 \\
GPT-6 Luna & sakaku's & 0.891 & 0.465 & 61 \\
GPT-6.1 Sol & tascj's & 0.928 & 0.492 & 63 \\
GPT-6.1 Sol & sakaku's & 0.932 & 0.350 & 44 \\
GPT-6 Astra & sakaku's & \textbf{0.936} & \textbf{0.329} & \textbf{35} \\
\bottomrule
\end{tabular}
\end{center}
Two lessons in one table: a stronger model helps, and \textbf{the instructions matter as much as the model}. sakaku's
instructions for turning evidence into a probability halved the confident errors.

\subsection*{Calibration: taming over-confidence}

\begin{figure}[H]
\centering
\begin{tikzpicture}
\begin{axis}[width=6.4cm, height=6.4cm, xmin=0, xmax=1, ymin=0, ymax=1, xlabel={agent's probability $p$},
    ylabel={calibrated label $p'$}, grid=major, xlabel style={font=\footnotesize}, ylabel style={font=\footnotesize}]
  \addplot[muted, dashed] coordinates {(0,0) (1,1)};
  \addplot[them, line width=1.4pt, domain=0.001:0.999, samples=120] {1/(1+exp(-0.77*ln(x/(1-x))))};
  \node[font=\scriptsize, anchor=west, text=them!85!black] at (axis cs:0.05,0.88) {$p' = \sigma(0.77\cdot\text{logit}\,p)$};
\end{axis}
\end{tikzpicture}\hspace{6mm}
\begin{minipage}[b]{0.42\linewidth}\small
\term{Temperature scaling}: convert the probability back to a logit, multiply by a factor $k$, squeeze again.
With $k = 0.77$ (fitted on gold), a 0.99 becomes about 0.97 and a 0.02 about 0.05: the agents were a bit too sure
of themselves. In tascj's step-by-step test this added +0.0014 on gold.\par\bigskip
\end{minipage}
\caption{Calibrating the labels before training.}
\end{figure}

\section{A guide that learns from the doctors}

\begin{figure}[H]
\centering
\begin{tikzpicture}[x=1mm,y=1mm,
  st/.style={draw=them, fill=whyfill, rounded corners=2mm, minimum width=30mm, minimum height=10mm, align=center, font=\scriptsize},
  ar/.style={-{Stealth[length=2mm]}, line width=0.8pt, them}]
  \node[st] (a) at (0,0) {label one gold study\\ as usual};
  \node[st] (b) at (40,0) {see the doctors'\\ answers};
  \node[st] (c) at (80,0) {reflect and rewrite\\ a notes file};
  \draw[ar] (a) -- (b); \draw[ar] (b) -- (c);
  \draw[ar] (c.south) -- ++(0,-6) -| (a.south) node[pos=0.25, below, font=\scriptsize] {next gold study};
  \node[st, draw=wordline, fill=wordfill, minimum width=36mm] (d) at (125,0) {after N studies:\\ write a \textbf{reading guide}};
  \draw[ar, wordline] (c) -- (d);
\end{tikzpicture}
\caption{The ``trained guide'': one long agent session learns from the 58 gold studies and writes down what it learned.
The guide is then added to the instructions for labelling everything else.}
\end{figure}

\begin{center}\small
\begin{tabular}{lccc}
\toprule
Astra labels on the 58 gold studies & AUC & log loss & confident errors \\
\midrule
without the guide & 0.936 & 0.329 & 35 \\
+ guide learned on other folds (honest test) & 0.936 & 0.307 & \textbf{19} \\
+ guide learned on the same 58 (optimistic) & 0.943 & 0.300 & 16 \\
\bottomrule
\end{tabular}
\end{center}
The honest test shows the real effect: same ranking, but \textbf{half the confident errors}.

\section{Small decisions that mattered}

\begin{itemize}
\item \textbf{Continuous, not binned.} Asking for one of six bins gave 0.883; asking for a probability 0.936.
      Our v4 labeller had the same instinct (probabilities over letters), but only five letters.
\item \textbf{Silence is not negative, uncertainty is not positive}: the same rule as our v4 letters N and U.
\item \textbf{Synovitis from the images.} Synovitis wording was removed from the report before the agent saw it, so the
      agent had to judge the joint lining from the images instead of copying the report.
\item \textbf{Mixing label sources.} The final labels for some models averaged five sources (two Astra versions, the
      teammate's Astra labels, two single-pass readers) with one weight per source for all findings. Per-finding
      weights looked better but did not survive a leave-one-out check.
\end{itemize}

\begin{keyidea}
The decisive difference was not a cleverer prompt. It was \textbf{giving the labeller the images}. Reports describe
what the original radiologist chose to write; the gold labels describe what is in the images. Only an image-reading
labeller can close that gap, and that is exactly where our labels were weakest (Synovitis, Contusion, Lateral OA).
\end{keyidea}

\begin{watchout}
This was within the rules: the host had explicitly allowed hosted LLM APIs for working with the reports. The cost was
real (thousands of agent sessions with a frontier model), but by the rules it counted as ``reasonably accessible''.
\end{watchout}

\begin{checkpoint}
\begin{itemize}
\done I can explain the difference between a prompt and an agent.
\done I know why evidence-first answers and calibration made better training labels.
\quiz Q12.1: Why might a label set with a higher AUC still be worse for training than one with a lower AUC?
\quiz Q12.2: Why remove the synovitis wording from the report before the agent reads it?
\quiz Q12.3: The trained guide was learned on gold studies. Why is the ``other folds'' row the one to trust?
\end{itemize}
\end{checkpoint}
